\section{Paper Survey}
\label{sec:paper_survey}

Recent research has examined the reliability of Multilingual Retrieval-Augmented Generation (mRAG) systems, pointing out clear vulnerabilities in how models process cross-lingual context.

\begin{enumerate}
    \item \textcite{wang2026} showed that mRAG rerankers suffer from strong language bias, often overlooking key documents that are not written in English or the query language. While their work focuses on upstream document retrieval and reranking, the proposed project focuses on the downstream generator, examining how the LLM resolves conflicting evidence after documents have already been retrieved.

    \item \textcite{ki2025} analyzed model internals to show that LLMs tend to cite English sources when given a mix of multilingual documents, sometimes favoring language over passage relevance. Whereas their study measures document citation frequency in general multi-document mixtures, the proposed work directly pits two mutually exclusive factual claims against each other to test which fact the model adopts as true.

    \item \textcite{park2025} found that mRAG generators prefer the query language or Latin scripts in natural web passages, leading to inconsistent answers across languages. In contrast to their use of unstructured web passages where length and style vary, the proposed project uses controlled minimal-pair sentence templates and a closed-book baseline to isolate language identity from passage style and prior knowledge.

    \item \textcite{qi2025} evaluated how models utilize multilingual context and found that while LLMs can extract facts from foreign languages, irrelevant distractor text in the query language significantly hurts performance. Rather than testing irrelevant distractors against relevant facts, the proposed work tests two equally relevant but directly contradictory claims in English and Korean.

    \item \textcite{li2025} studied \enquote{language drift} in mRAG, showing that English acts as a latent semantic attractor that pulls the model toward generating English tokens during decoding. While their work focuses on the output language of the generated text, the proposed project investigates whether English also acts as an epistemic attractor that determines which factual claim the model trusts.
\end{enumerate}